%% MeshSOS --- IEEE conference paper
%% Compile on Overleaf: New Project -> Upload Project -> this file,
%% or paste into Overleaf's "IEEE Conference Template". IEEEtran is
%% preinstalled there. Local pdflatex needs texlive-publishers.

\documentclass[conference]{IEEEtran}
\IEEEoverridecommandlockouts
\usepackage{cite}
\usepackage{amsmath,amssymb,amsfonts}
\usepackage{algorithmic}
\usepackage{graphicx}
\usepackage{textcomp}
\usepackage{xcolor}
\usepackage{booktabs}
\usepackage{url}
\def\BibTeX{{\rm B\kern-.05em{\sc i\kern-.025em b}\kern-.08em
    T\kern-.1667em\lower.7ex\hbox{E}\kern-.125emX}}

\begin{document}

\title{Fitting an Authenticated Emergency Report into One BLE Advertising Frame:
Semantic Compression for Infrastructure-Free Disaster Mesh Networks}

\author{
\IEEEauthorblockN{Gauri Dhingra}
\IEEEauthorblockA{\textit{Dept. of CSE (AI\&ML)} \\
\textit{KIET Deemed to be University}\\
Ghaziabad, India \\
gaurdhingra0@gmail.com}
\and
\IEEEauthorblockN{Ity Gaur}
\IEEEauthorblockA{\textit{Dept. of CSE (AI\&ML)} \\
\textit{KIET Deemed to be University}\\
Ghaziabad, India \\
itygaur@gmail.com}
\and
\IEEEauthorblockN{Mahima Bhat}
\IEEEauthorblockA{\textit{Dept. of CSE (AI\&ML)} \\
\textit{KIET Deemed to be University}\\
Ghaziabad, India \\
mahimabhat2@gmail.com}
\and
\IEEEauthorblockN{Nischhal Garg}
\IEEEauthorblockA{\textit{Dept. of CSE (AI\&ML)} \\
\textit{KIET Deemed to be University}\\
Ghaziabad, India \\
nischhalgarg61@gmail.com}
\and
\IEEEauthorblockN{Preeti Singh}
\IEEEauthorblockA{\textit{Dept. of CSE (AI\&ML)} \\
\textit{KIET Deemed to be University}\\
Ghaziabad, India \\
preetisinghthakur07@gmail.com}
}

\maketitle

\begin{abstract}
When cellular infrastructure fails, smartphone mesh networks can still carry
distress reports hop by hop over Bluetooth Low Energy. On the handsets most
common in disaster-prone regions the only universally supported channel is
legacy BLE advertising, which carries 31 bytes per frame and about 27 after
the mandatory AD structures. An authenticated emergency report does not fit:
in our system it occupies 124 bytes before a single character of free text,
because an Ed25519 public key and signature account for 96 of them. We present
a semantic compression path that fits a complete authenticated report into one
frame. Free text is encoded on the handset into a 3-byte envelope carrying
category, casualty count, mobility, hazard flags, injury severity and the
encoder's own confidence; the message is then packed into 25 bytes alongside a
coarse grid location and a 4-byte authentication tag. Reaching that budget
required reducing cryptographic overhead as well as text, and we state each
trade explicitly. Below a confidence threshold nothing is compressed and the
original text is transmitted. We extend a discrete-event mesh simulator to
model frame-level fragmentation, size-dependent airtime and receiver-side
collisions, none of which it previously represented. Across 60 trials per
configuration, compression raises the fraction of reports reaching a gateway
from 53.3\% to 60.0\% at 200 nodes and cuts channel occupancy by 8.5--15.7\%.
We also report two negative results: a model trained on synthetic data reaches
99.7\% on its own held-out split and loses to a keyword baseline on unseen
text, though a cascade between them reaches 67.7\% against 59.5\% and 53.4\%
alone; and a sample of 520 tweets from a standard crisis corpus proved 51\%
commentary and 18\% news reporting, a mismatch of population rather than of
annotation quality.
\end{abstract}

\begin{IEEEkeywords}
disaster communication, BLE mesh, semantic compression, on-device NLP,
emergency triage, code-mixed text
\end{IEEEkeywords}

\section{Introduction}

Earthquakes, floods and building collapses remove the infrastructure that
emergency communication depends on. Base stations lose power, backhaul is cut,
and the handsets that survive have no network to reach. Smartphone mesh
networking offers a way through: phones relay for one another over short-range
radio until a message reaches a device that still has connectivity.

The constraint that shapes such a system is not bandwidth in the usual sense.
It is the advertising frame. BLE connectionless advertising is the only mode
that works without pairing, without a connection, and on every Android handset
in the field, including the five-year-old devices that dominate the regions
where these systems would be deployed. A legacy advertising PDU carries 31
bytes; after the mandatory Flags structure and a 16-bit Service Data header,
roughly 27 remain.

Our heartbeat message was engineered to 25 bytes and fits. The SOS message does
not. Measured on our implementation, an SOS occupies 124 bytes carrying no text
at all, rising to 162--187 bytes with a short description --- six or seven
frames. The reason is not the text. An Ed25519 public key (32\,B) and signature
(64\,B) account for 96 of those 124 bytes, 77\% of the message, before anything
is said about the emergency.

This matters because fragmentation is not free. A message split across $n$
frames arrives only if every frame arrives: at 3\% per-frame loss, a six-frame
message fails 16.7\% of the time where a single-frame message fails 3\%. It
also occupies the channel longer, and in a dense mesh where every node relays,
channel occupancy limits how many reports get through.

Fig.~\ref{fig:pipeline} shows the path a report takes. The question this paper asks is therefore narrower than whether NLP can help
disaster response. It is: \emph{can a complete, authenticated emergency report
be made to fit in one legacy BLE advertising frame, and what does that buy at
the network layer?}

\subsection{Contributions}

\textbf{A semantic envelope sized against a derived budget.} We compress
free-form emergency text into 3 bytes carrying six fields and fit the complete
authenticated message into 25. The budget follows from the advertising PDU
structure rather than being asserted, and reaching it required reducing
cryptographic overhead as well as text --- a trade prior work has not had to
confront because it does not target a frame boundary.

\textbf{A size-aware mesh simulator, and network-level results.} We add
frame-level fragmentation, size-dependent airtime and receiver-side collisions
to a discrete-event simulator that previously charged every transmission a
fixed 5\,ms and dropped whole messages with one coin flip, so payload length
could not affect any outcome. We then measure what compression does to
delivery, not only to classification accuracy.

\textbf{Negative results reported in full.} A model trained on synthetic data
outperforms a keyword baseline by nearly ten times on its own held-out split
and loses to it on hand-written text. Stratifying 10{,}000 messages by distance
from the lexicon shows the two fail in opposite places, and a cascade between
them beats either. We also document that a standard crisis corpus does not
contain the kind of message this system carries.

\section{Related Work}

\subsection{Classification of crisis text}
Automatic categorisation of disaster-related messages is mature. HumAID
\cite{humaid} and the consolidated CrisisBench benchmark \cite{crisisbench}
provide standard splits and transformer baselines for humanitarian category,
informativeness and priority labelling. Subsequent work has largely refined
accuracy on these benchmarks rather than altering what is transmitted. We treat
crisis-text classification as available infrastructure rather than as a
contribution, and adopt it only as the front end of an encoder whose output is
a transmission format rather than a label.

\subsection{On-device intelligence in offline mesh networks}
The closest work is ResQMesh \cite{resqmesh}, which extends an open-source BLE
mesh messenger with on-device machine learning. Its priority module combines
roughly ninety keyword patterns with a quantised INT8 Conv1D TensorFlow Lite
classifier, assigns nine categories and four urgency tiers, and reports a
composite macro-F1 of 0.920 on 405 held-out messages at 0.19--1.97\,ms.
Classified messages enter a priority queue so critical traffic preempts other
traffic at the radio layer.

ResQMesh establishes that useful semantic processing is feasible on commodity
Android hardware without connectivity, and we build on that premise.
\textbf{Our work differs in what the classifier is for.} ResQMesh uses the
inferred category to \emph{order} transmissions; the message itself is
unchanged, and no payload size, MTU, advertising budget or fragmentation
behaviour is reported. A message too large for one frame remains too large
after classification, and the reported latency improvement is measured at the
queue rather than over the air. We use the inferred semantics to
\emph{constrain the payload} and evaluate the consequence at the network layer.

\subsection{Semantic communication and extreme compression}
Semantic communication reframes transmission around conveying meaning rather
than reconstructing symbols. Recent work examines tiny language models under a
joint latency--accuracy--size constraint \cite{tlm6g}, and semantic text
compression evaluated by preserved downstream classification rather than
reconstruction fidelity \cite{semcomp}. This literature assumes a channel with
meaningful capacity; to our knowledge none targets a single-digit-byte residual
budget inside a 31-byte advertising PDU. In the grey literature, RelayZero
\cite{relayzero} proposes offline edge-AI semantic compression for disaster
mesh using a locally hosted language model. We note it as concurrent evidence
of interest and distinguish our contribution: it reports no accuracy, latency
or reliability measurement, its mesh layer is a local TCP socket rather than a
radio, its stated 25-byte payload is inconsistent with its own published
example string, and its encoder is not deployable on the handsets we target.

\subsection{Authenticated offline emergency broadcast}
The Offline Emergency Peer-to-Peer Broadcast Protocol (OEPB) \cite{oepb}, an
experimental Internet-Draft, specifies authenticated emergency alerts over
short-range radios: a 256-byte CBOR datagram with a 40-byte fixed header and an
optional 64-byte Ed25519 signature. Its existence confirms independently that
authenticated infrastructure-free emergency broadcast is worth specifying.
\textbf{Its answer is the one we set out to avoid.} On BLE 4.x, OEPB states
that the 31-byte advertisement payload requires L2 fragmentation and
reassembly; its signed overhead alone is 104 bytes, and free text is carried as
plain UTF-8 under a byte cap with no semantic compression. We ask whether that
fragmentation can be avoided rather than accommodated.

\subsection{Short authentication tags on constrained links}
Our 4-byte tag has prior art. Cumulative Message Authentication Codes
\cite{cumac} address the same tension: a full MAC does not fit, so a 128-bit
tag is split into segments spread across successive messages, with the receiver
accumulating strength over time. We adopt the truncated tag without the
accumulation, and note the omission as the clearest available improvement: our
store-and-forward layer already re-broadcasts each SOS every ten seconds for up
to six hours, so the repeated transmissions such a scheme requires are present
and currently carry an identical tag each time.

\subsection{Inference from absence of signal}
Post-disaster victim localisation is typically a sensing problem, addressed
with millimetre-wave radar, UAV-mounted sensors, or passive detection of
emissions from victims' handsets \cite{victim}. These detect a victim's
presence. The silence-inference component of the system this work extends
inverts the premise: participants carry a cooperating application emitting
periodic heartbeats, and a disaster is inferred from the \emph{correlated
cessation} of those heartbeats within a spatial cluster, discriminated against
battery depletion, deliberate shutdown, a single participant walking out of
range, and whole-network partition.

\subsection{Code-mixed emergency text}
Our deployment context produces Hindi--English code-mixed input. Pretrained
resources exist, including L3Cube-HingCorpus and HingBERT \cite{hingbert}, but
have not to our knowledge been applied to emergency triage under a
transmission-size constraint.

\section{Background: The Byte Budget}

\subsection{Why legacy advertising}
Bluetooth~5 extended advertising carries up to 254 bytes per packet and 1{,}650
via chaining, but requires a Bluetooth 5.0 chipset and OEM driver support.
Budget and older Android handsets --- precisely those most common in
disaster-prone regions --- support only legacy advertising. Wi-Fi Direct and
Nearby Connections offer higher throughput but are likewise not universally
available and cost substantially more energy. A system whose emergency path
depends on any of these degrades to nothing on the devices most likely to be
present. We therefore treat legacy advertising as the floor every message must
clear, and faster transports as an optimisation when available.

\subsection{Deriving the usable payload}

\begin{table}[htbp]
\caption{Usable payload of a legacy advertising PDU}
\begin{center}
\begin{tabular}{lr}
\toprule
\textbf{Element} & \textbf{Bytes} \\
\midrule
Advertising PDU payload & 31 \\
Flags AD structure (\texttt{0x02 0x01 0x06}) & $-3$ \\
Service Data AD header (length, type, 16-bit UUID) & $-4$ \\
\midrule
\textbf{Usable for application data} & \textbf{27} \\
\bottomrule
\end{tabular}
\label{tab:budget}
\end{center}
\end{table}

\subsection{What the existing messages cost}

\begin{table}[htbp]
\caption{Message sizes before compression}
\begin{center}
\begin{tabular}{lrr}
\toprule
\textbf{Message} & \textbf{Bytes} & \textbf{Frames at 27\,B} \\
\midrule
Heartbeat & 25 & 1 \\
SOS, no text & 124 & 5 \\
SOS, 13\,B text & 137 & 6 \\
SOS, 37\,B text & 161 & 6 \\
SOS, 200\,B text & 324 & 12 \\
\bottomrule
\end{tabular}
\label{tab:sizes}
\end{center}
\end{table}

\begin{figure}[htbp]
\centerline{\includegraphics[width=\columnwidth]{figures/fig_budget.pdf}}
\caption{Message size against the frame budget. Everything above the dashed
line fragments; 96 of the 124-byte floor is the public key and signature, so
compressing text alone cannot cross it.}
\label{fig:budget}
\end{figure}

The heartbeat was designed to the budget and fits. The SOS does not, and the
reason is cryptography rather than content: a 13-byte header, a 15-byte body, a
32-byte public key and a 64-byte signature give a 124-byte floor.
Fig.~\ref{fig:budget} shows this against the budget. \textbf{Compressing text alone moves a typical report from six frames to five.}
It cannot reach one.

\subsection{Why frame count matters}
\textbf{Compound loss.} A fragmented message arrives only if every fragment
arrives. With independent per-frame loss $p$ and $n$ frames, delivery
probability is $(1-p)^n$. At $p = 0.03$, one frame fails 3\% of the time and
six frames fail 16.7\% --- a $5.6\times$ increase in failure for the same
report.

\textbf{Channel occupancy.} On the 1M PHY each bit takes 1\,\textmu s, and each
frame carries 10 bytes of PDU overhead plus an inter-frame space. A 25-byte
message occupies about 0.43\,ms; a 364-byte message across 14 frames occupies
about 6.1\,ms. In a flooding mesh this difference multiplies across the network
and raises the chance that two frames overlap at a receiver.

\begin{figure*}[t]
\centerline{\includegraphics[width=\textwidth]{figures/fig_pipeline.pdf}}
\caption{The compression path. A report is encoded on the handset into three
bytes and sent as a 25-byte authenticated frame; the gateway expands it for
triage. When the encoder's confidence falls below threshold nothing is
inferred and the original text is sent instead, at six or seven frames.}
\label{fig:pipeline}
\end{figure*}

\section{System Design}

\subsection{The semantic envelope}
Free text is encoded into 24 bits: category (4), casualty count (3, where 7
means seven or more), mobility (2), hazard flags (6), severity (2), the
encoder's own confidence (3), and 4 reserved bits. Confidence travels
\emph{with the data} deliberately: a rescuer must see how much of a report was
inferred rather than stated, and the relay layer uses the same value to decide
whether to compress at all.

\subsection{Reaching 27 bytes}
Text compression alone is insufficient. Three further reductions were
necessary, each a trade we state rather than hide. The public key is resolved
from \texttt{node\_id} at the gateway rather than transmitted ($-32$\,B), so
only enrolled nodes can send compact messages. The 64-byte Ed25519 signature
becomes a 4-byte truncated HMAC ($-60$\,B); this is \emph{not} a digital
signature, it authenticates only to a gateway holding the shared key, and a
32-bit tag admits a one-in-four-billion forgery per attempt. Float latitude and
longitude become the ${\sim}100$\,m grid cell the heartbeat already uses
($-4$\,B), giving rescuers a cell rather than a pin, with ${\sim}11$\,m
quantisation error observed.

\begin{table}[htbp]
\caption{The compact SOS message}
\begin{center}
\begin{tabular}{lr}
\toprule
\textbf{Element} & \textbf{Bytes} \\
\midrule
Header (version, type, node id, seq, TTL, timestamp) & 13 \\
Grid $x$, grid $y$ & 4 \\
Battery & 1 \\
Semantic envelope & 3 \\
Authentication tag & 4 \\
\midrule
\textbf{Total} & \textbf{25} \\
\bottomrule
\end{tabular}
\label{tab:compact}
\end{center}
\end{table}

\textbf{The tag excludes TTL.} TTL decrements at every hop, so a tag computed
over the live header would verify only at the originating node. Zeroing it
means the tag authenticates the report rather than the path it travelled. The
grid is relative to a deployment origin distributed at enrolment: absolute
coordinates do not fit, since longitude 77.209 at $0.001^\circ$ resolution is
77{,}209, past the 32{,}767 ceiling of a signed 16-bit field. Both SOS layouts
share the message type; the version byte distinguishes them.

\subsection{Threat model}
Four bytes of authentication invites an immediate objection, so we state
precisely what the tag defends, against whom, and what it knowingly does not.
The asset is the integrity of the rescuer's view: a forged SOS moves rescuers
to an empty field. Confidentiality is deliberately absent --- an encrypted cry
for help trades availability in exactly the wrong place --- so anyone in radio
range learns that grid cell $(x,y)$ reported trouble; we state this privacy
cost rather than hide it.

The adversary transmits and receives freely but holds no key. Blind forgery
succeeds per attempt with probability $2^{-32}$; under the airtime model of
Section~III a 25-byte frame occupies 430\,\textmu s, bounding one channel at
${\sim}2{,}300$ attempts/s, so even an attacker saturating the channel --- a
visible denial of service in itself --- expects a first forgery after
$2^{32}/2300 \approx 21$ days of continuous transmission, and the gateway can
rate-limit verification failures per node id. Replay is bounded because the
tag covers \texttt{node\_id}, \texttt{seq} and \texttt{timestamp} (only TTL
is zeroed), and both relays and gateway deduplicate on
(\texttt{node\_id}, \texttt{seq}).

Two honest weaknesses remain. The prototype uses a single group key, so one
captured device forges any identity; the deployment design derives
$K_{node} = \mathrm{KDF}(K_{master}, \mathit{node\_id})$ at enrolment, at
zero byte cost, after which a compromise forges only itself. And relays,
holding no key, amplify injected frames with fresh sequence numbers;
suppression ($k{=}2$) and TTL bound the amplification, a per-node token
bucket at relays is designed but unimplemented, and an attacker with a
transmitter could in any case jam the band outright. Tag comparison is
constant-time at both endpoints. The clearest upgrade is CuMAC-style
accumulation (Section~II): rotate four segments of a 16-byte tag across the
store-and-forward re-broadcasts the system already performs, reaching 128-bit
strength over repeats at zero bytes per frame.

\subsection{The encoder}
The deployed encoder is rule-based: lexicon matching over category, casualty,
mobility, hazard and severity terms, with numeric extraction that binds a count
to a nearby person-word so ``4 people trapped'' is read and ``4 floors'' is
not. Lexicons cover English and romanised Hindi together, since code-mixed
input is the norm rather than an edge case. Negation is checked on both sides
of a match, because English negates before the term (``no injuries'') and
romanised Hindi after it (``\emph{ghayal nahi}''). Confidence combines how many
of the five fields were recovered with how much of the message the lexicons
failed to account for: \textbf{a message carrying much the encoder could not
read lowers confidence rather than raising it}.

\subsection{Fallback}
Below a confidence threshold of 0.45 the sender transmits the original
full-text SOS. Nothing is inferred that the encoder did not understand.

\section{Evaluation}

\subsection{Making the simulator size-aware}
The simulator this work extends could not have measured compression. Every
transmission was charged a fixed 5\,ms regardless of length, whole messages
were dropped with a single coin flip, and there was no collision model. We
verified the consequence: the same 200-node scenario run with an 11-byte and a
240-byte message produced byte-identical results --- 0 delivered, 2{,}474
transmissions, 5{,}588 dedupe drops in both. We added airtime derived from
length, fragmentation above 27 bytes with independent per-frame loss, and
collisions between frames overlapping at one receiver, each switchable so prior
behaviour remains reproducible.

\textbf{A note on a prior published figure.} Our earlier evaluation reported 0
SOS delivered at 200 nodes. That is a single-seed artefact: seed 42 places the
source 17 hops from the nearest gateway against a TTL of 8, while the median
across seeds is 6. All results below average over seeds.

\subsection{Encoder fidelity, stratified}
One number over a generated set would report how the generator happened to
phrase things. We therefore split 10{,}000 messages by how far the wording sits
from the encoder's lexicon, with labels true by construction.

\begin{table}[htbp]
\caption{Rule encoder by stratum (2{,}000 messages each)}
\begin{center}
\begin{tabular}{lrrrr}
\toprule
\textbf{Stratum} & \textbf{Category} & \textbf{Count} & \textbf{All} & \textbf{Fallback} \\
\midrule
known\_vocab      & 79.8\% & 100.0\% & 58.6\% & 16.2\% \\
synonyms (OOV)    &  8.3\% &  60.3\% &  0.1\% & 97.0\% \\
spelling          & 69.8\% & 100.0\% & 50.9\% & 25.2\% \\
code\_mixed       & 79.9\% & 100.0\% & 60.4\% & 17.4\% \\
unintelligible    & ---    & ---     & ---    & 100.0\% \\
\bottomrule
\end{tabular}
\label{tab:strata}
\end{center}
\end{table}

\textbf{The out-of-vocabulary cliff is total, and it is safe.} On synonyms the
encoder recovers the category 8.3\% of the time, and refuses to compress 97.0\%
of them. The failure mode is slower transmission, not a wrong report.
Code-mixed input is not the weak point: at 60.4\% it edges out pure
known-vocabulary text. Refusal on unintelligible input is perfect across 2{,}000
messages.

\subsection{A learned encoder}
The real corpora were unreachable from our environment, so we generated 6{,}000
messages from templates whose field values are known by construction --- ground
truth rather than rule output --- and trained a character n-gram TF-IDF model
with one logistic-regression head per field. It reaches 99.7\% exact-envelope
accuracy on held-out synthetic data against the rules' 10.3\%, and 28.0\%
against the rules' 76.0\% on hand-written messages, at 808\,KB and eight times
the latency. The model learned the generator, not the task.

\subsection{The two encoders fail in opposite places}

\begin{table}[htbp]
\caption{Category accuracy: rules, model, and a cascade between them}
\begin{center}
\begin{tabular}{lrrr}
\toprule
\textbf{Stratum} & \textbf{Rules} & \textbf{Neural} & \textbf{Cascade} \\
\midrule
known\_vocab    & 79.8\% & 62.3\% & \textbf{81.2\%} \\
synonyms (OOV)  &  8.3\% & \textbf{39.1\%} & 38.8\% \\
spelling        & 69.8\% & 61.7\% & \textbf{74.8\%} \\
code\_mixed     & \textbf{79.9\%} & 50.5\% & 75.8\% \\
\midrule
\textbf{Overall} & 59.5\% & 53.4\% & \textbf{67.7\%} \\
\bottomrule
\end{tabular}
\label{tab:cascade}
\end{center}
\end{table}

\begin{figure}[htbp]
\centerline{\includegraphics[width=\columnwidth]{figures/fig_cascade.pdf}}
\caption{The two encoders fail in opposite places. The lexicon wins where it
has the words; the model wins on wording it has never seen. Routing between
them on the confidence the rule encoder already emits beats either alone.}
\label{fig:cascade}
\end{figure}

Scored per stratum, the model does not simply lose: on out-of-vocabulary
wording it recovers nearly five times what the lexicon does, because character
n-grams reach words no keyword list contains. The rule encoder already emits a
confidence deciding whether to compress; routing on it requires nothing further
to be trained. The unintelligible stratum is excluded from the overall figure,
since its gold label is OTHER and both encoders answer OTHER. Unintelligible
input is still refused 100\% of the time.

\subsection{Why not general-purpose compression}

\begin{table}[htbp]
\caption{Compressing the text instead of encoding it. Methods are scored on
3{,}120 held-out real messages; the shared dictionary is built from a disjoint
7{,}278-message training split. \emph{Fits} asks whether the resulting SOS
occupies one frame: for the first three it never can, because 124 bytes of
header and cryptography precede the text.}
\begin{center}
\begin{tabular}{lrrl}
\toprule
\textbf{Method} & \textbf{Mean\,B} & \textbf{Max} & \textbf{Fits 27\,B?} \\
\midrule
Raw UTF-8                 &  91.7 & 585 & no \\
zlib (level 9)            &  75.2 & 341 & no \\
zlib + shared dictionary  &  59.3 & 326 & no \\
Semantic envelope         & \textbf{3.0} & \textbf{3} & \textbf{always} \\
\bottomrule
\end{tabular}
\label{tab:compress}
\end{center}
\end{table}

General-purpose compressors \emph{enlarge} messages this short: their headers
exceed what they save. The serious competitor is zlib with a shared
dictionary, which is lossless and genuinely helps --- on held-out real
messages it cuts the text to 65\% of raw. We distinguish it on three grounds.
Its output is variable, and a frame budget is a bound rather than an average.
It yields bytes rather than fields, so the dashboard must still parse text to
rank by urgency. And it requires the dictionary to stay synchronised: a
mismatch destroys the message rather than degrading it.

\textbf{A correction.} An earlier draft of this table reported 11.2 bytes for
the shared dictionary, measured on 25 hand-written messages. That figure does
not survive a held-out split. Rebuilding the dictionary from a disjoint
training corpus gives 36.9 bytes on those same 25 messages and 59.3 bytes
across 3{,}120 held-out real ones. A dictionary that contains the evaluation
text makes each message a literal substring of it, and zlib then emits little
more than a back-reference: we measure 4.6 bytes in that condition. Any
shared-dictionary figure in the single digits should be read as a split that
leaked, including ours.

The correction does not change the conclusion, and sharpens it. Compressing
the fallback text from 91.7 to 59.3 bytes moves a full SOS from 8.46 to 7.26
advertising frames --- a real 14\% saving we intend to ship --- but the fixed
part of that message is 124 bytes, five frames, before a single character of
text, of which 96 bytes are the signature and public key. \emph{Even
compressing the text to nothing leaves five frames against the compact path's
one.} The bound is the cryptography, not the language.

\subsection{Network effect}
Three conditions, \textbf{500 paired trials per cell} pairing seed and message
round-robin, TTL~8, 3\% per-frame loss: \textbf{A} full text, \textbf{B}
always compressed, \textbf{C} compressed when confident. All three conditions
see the same topology and the same message in every trial, so the design is
paired and the arms are compared with an exact McNemar test on the discordant
trials. Intervals are 95\% Wilson score intervals, which stay inside $[0,1]$
at the rates observed here where the normal approximation does not.

\begin{table}[htbp]
\caption{SOS delivered to a gateway, 500 paired trials per cell}
\begin{center}
\begin{tabular}{lrrrcr}
\toprule
\textbf{Nodes} & \textbf{A} & \textbf{B} & \textbf{C} & \textbf{95\% CI (C)} & \textbf{Airtime (C)} \\
\midrule
50  & 84.6\% & 91.2\% & 88.8\% & [85.7, 91.3] & 13.7\% \\
100  & 72.8\% & 82.6\% & 79.8\% & [76.1, 83.1] & 10.6\% \\
200  & 54.8\% & 63.8\% & 61.2\% & [56.9, 65.4] & 7.2\% \\
\bottomrule
\end{tabular}
\label{tab:network}
\end{center}
\end{table}

\begin{figure}[htbp]
\centerline{\includegraphics[width=\columnwidth]{figures/fig_delivery.pdf}}
\caption{SOS reports reaching a gateway, by condition and network size, with
95\% Wilson intervals. The margin widens as the mesh loads: compression buys
the most where contention is highest.}
\label{fig:delivery}
\end{figure}

Every A-versus-C difference is significant at $p < 0.001$. At 200 nodes the
paired counts are 34 trials where C delivered and A did not against 2 the
other way; at 100 nodes, 37 against 2; at 50 nodes, 22 against 1. The relative
gain grows with load --- $+5.0\%$, $+9.6\%$, $+11.7\%$ at 50, 100 and 200
nodes --- which is the expected direction if the mechanism is contention
rather than anything particular to the payload.

An earlier version of this evaluation used 60 trials per cell. At that size
the 95\% interval on a rate near 0.8 is roughly $\pm 10$ percentage points,
wide enough to contain every effect reported here, and the apparent ordering
of the arms moved between runs. The numbers above supersede it.

\textbf{C is the shipping configuration and the number to quote}; B assumes
the encoder is never unsure, which is false for 28\% of these messages. B
beats C significantly at every size ($p < 0.002$), and that gap is precisely
the cost of the fallback --- an honest encoder that refuses to guess gives up
part of the gain, and we report the part it gives up rather than quoting the
ceiling. A separate sweep separates TTL-limited from size-limited failure:
where TTL is too low for the topology, both short and long messages fail and
compression cannot help.

\section{Discussion}

\subsection{What the negative result means}
Template-generated text teaches a model the template. On its own split that
model looks close to perfect; on text written by a person it collapses below a
keyword matcher. Crisis-NLP systems are frequently evaluated on held-out splits
of the corpus they were trained on, and where that corpus is narrow the
reported accuracy may say more about the corpus than the system. Our result is
an extreme case because our corpus is synthetic, but the direction of the error
is the same one a single-event benchmark would produce.

\subsection{Available corpora collect the wrong half of the conversation}
We attempted to replace generated text with real data and did not succeed.
HumAID \cite{humaid} supplies 76{,}484 annotated tweets from 19 disasters. We
drew 520, weighted toward the classes most likely to contain an emergency, and
classified what we had: 51\% commentary, 18\% news reporting, 17\% retweets,
1\% donation appeals. The 13\% matching first-person patterns were largely
expressions of sympathy rather than requests for help.

This is not a defect in HumAID. It is a mismatch of population. Crisis corpora
are collected from public social media, where the people writing are
overwhelmingly \emph{observers}. A mesh system of this kind carries messages
from \emph{participants}. The two differ in length, register, person, and the
presence of hashtags and audience, and the labels of the first describe
information type rather than emergency type. The Figure Eight \emph{Disaster
Response Messages} corpus \cite{drm} carries a \texttt{direct} genre of
messages and SMS sent during the 2010 Haiti and Chile earthquakes, a closer
population, though translated and smaller.

\subsection{The safety case}
A compression error here is not a metric: reading ``two children trapped'' as
``one adult trapped'' sends the wrong team. Three bits of the envelope carry
how certain the encoder was, so a rescuer's display can mark an inferred report
as inferred. A message where the lexicons account for little of the text scores
\emph{lower}, not higher --- the encoder treats its own incomprehension as
evidence against compressing. Below threshold the original text is transmitted
at the cost of more frames, a deliberate exchange of reliability for fidelity
in the direction that fails safe. A regression test asserts that casualty
counts are never reported lower than ground truth. We make no claim that these
eliminate the risk.

\section{Limitations}

\textbf{No real evaluation corpus.} Fidelity is measured on generated text and
25 hand-written messages. Our attempt to use a real corpus found a population
mismatch (\S\,VI-B). \textbf{No hardware validation.} Everything is simulated;
the encoder has not run on a handset and no energy measurement exists.
\textbf{A simplified radio model.} No multipath, capture effect, or
hardware-specific advertising behaviour. \textbf{Weaker authentication.} The
4-byte tag is not a digital signature and the shared key in our prototype is a
placeholder that must not ship. \textbf{Coarser location.} A ${\sim}100$\,m
cell rather than a GPS pin. \textbf{Enrolment required.} An unenrolled handset
cannot send a compact SOS. \textbf{Single implementation.} Results have not
been reproduced against an independent BLE stack.

\section{Conclusion}

An authenticated emergency report can be made to fit in one legacy BLE
advertising frame, but not by compressing text alone --- cryptographic overhead
dominates the message, and reaching 27 bytes required reducing it too. Doing so
raises the share of reports reaching a gateway from 53.3\% to 60.0\% at 200
nodes and cuts channel occupancy by 8.5--15.7\%, with a confidence-gated
fallback that transmits the original text on the 23\% of messages the encoder
does not understand.

The encoder that achieves this is a keyword matcher. A learned model trained on
synthetic data scored 99.7\% on its own held-out split and 28.0\% on
hand-written text. But the two fail in opposite places, and a cascade routed on
a confidence signal that already existed reaches 67.7\% against 59.5\% and
53.4\% alone. We report both negative results in full because quoting the
synthetic column would have been a false finding, and because together they say
what the next step needs: not a better architecture, but emergency text written
by people in emergencies, with the envelope fields annotated. We looked for it
and found that the standard corpora collect the other half of the conversation.

\section*{Artefact Availability}
All code, measured results and reproduction commands are in the project
repository.

\begin{thebibliography}{00}
\bibitem{humaid} F. Alam, U. Qazi, M. Imran, and F. Ofli, ``HumAID:
Human-annotated disaster incidents data from Twitter with deep learning
benchmarks,'' in \emph{Proc. ICWSM}, 2021.
\bibitem{crisisbench} F. Alam, H. Sajjad, M. Imran, and F. Ofli,
``CrisisBench: Benchmarking crisis-related social media datasets for
humanitarian information processing,'' \emph{arXiv:2004.06774}, 2020.
\bibitem{resqmesh} O. Pliekhov, ``AI-driven message prioritization for offline
BLE mesh networks in disaster response scenarios,'' SSRN preprint,
doi:10.2139/ssrn.6428398, 2026.
\bibitem{tlm6g} ``Bridging the semantic gap in 6G: Tiny language models under
the latency--accuracy--size trilemma,'' preprint, 2026.
\bibitem{semcomp} ``Semantic text compression for classification,'' preprint,
2023.
\bibitem{relayzero} ``RelayZero: Offline edge-AI semantic compression for
disaster mesh networks,'' online technical report, Oct. 2026.
\bibitem{oepb} K. Sharma, ``Offline Emergency Peer-to-Peer Broadcast Protocol
(OEPB),'' Internet-Draft draft-sharma-oepb-00, experimental, Jun. 2026.
\bibitem{cumac} ``Cumulative message authentication codes for
resource-constrained IoT networks,'' \emph{arXiv:2001.05211}, 2020.
\bibitem{drm} Figure Eight / Appen, ``Disaster Response Messages,'' dataset.
\bibitem{victim} ``Victim detection and localization in emergencies,''
\emph{Sensors}, vol.~22, no.~21, p.~8433, 2022.
\bibitem{hingbert} R. Nayak and R. Joshi, ``L3Cube-HingCorpus and HingBERT: A
code mixed Hindi-English dataset and BERT language models,''
\emph{arXiv:2204.08398}, 2022.
\bibitem{btspec} Bluetooth SIG, \emph{Bluetooth Core Specification}, Generic
Access Profile, AD structure definitions.
\bibitem{androidble} Android Developers, ``Bluetooth Low Energy scanning and
background execution limits.''
\end{thebibliography}

\end{document}
